% !TEX program = xelatex
% StoryLens — Bộ câu hỏi & trả lời VẤN ĐÁP dành riêng cho Đào Sỹ Duy Minh (PM & AI/ML)
% Biên dịch: xelatex QA_DuyMinh.tex   (chạy 2 lần để có mục lục)
\documentclass[11pt]{article}

\usepackage{fontspec}
\setmainfont{Times New Roman}
\newfontfamily\headfont{Arial}
\usepackage[a4paper,margin=2cm]{geometry}
\usepackage{xcolor}
\usepackage[most]{tcolorbox}
\usepackage{titlesec}
\usepackage{enumitem}
\usepackage{fancyhdr}
\usepackage{amssymb}
\usepackage{array}
\usepackage{booktabs}
\usepackage[hidelinks]{hyperref}

\definecolor{crimson}{HTML}{C0271A}
\definecolor{ink}{HTML}{1A1712}
\definecolor{gold}{HTML}{9C6B10}
\definecolor{steel}{HTML}{285B7A}
\definecolor{soft}{HTML}{6E665D}
\definecolor{leaf}{HTML}{2F6B3A}

\newcommand{\arr}{\,\ensuremath{\rightarrow}\,}
\newcommand{\lte}{\ensuremath{\leq}}
\newcommand{\gte}{\ensuremath{\geq}}
\newcommand{\x}{\ensuremath{\times}}
\newcommand{\code}[1]{\texttt{\small #1}}

% Nhãn mức độ / loại câu
\newcommand{\tagCore}{\textcolor{steel}{\footnotesize\textbf{[Cốt lõi]}}}
\newcommand{\tagTrap}{\textcolor{crimson}{\footnotesize\textbf{[BẪY]}}}
\newcommand{\tagHonest}{\textcolor{gold}{\footnotesize\textbf{[Trung thực]}}}

\titleformat{\section}{\Large\bfseries\headfont\color{crimson}}{\thesection.}{0.5em}{}
\titleformat{\subsection}{\large\bfseries\headfont\color{ink}}{}{0em}{}
\titlespacing{\section}{0pt}{14pt}{6pt}

\pagestyle{fancy}\fancyhf{}
\lhead{\footnotesize\headfont\textcolor{soft}{StoryLens · Ôn vấn đáp — Đào Sỹ Duy Minh}}
\rhead{\footnotesize\headfont\textcolor{soft}{\thepage}}
\renewcommand{\headrulewidth}{0.4pt}

% Khối câu hỏi–trả lời
\newtcolorbox{qa}[1]{breakable,enhanced,sharp corners,
  boxrule=0.8pt,colframe=crimson,colback=crimson!3,
  coltitle=white,colbacktitle=crimson,fonttitle=\bfseries\headfont\small,
  title={#1},left=7pt,right=7pt,top=5pt,bottom=6pt,
  before skip=7pt,after skip=7pt}
\newcommand{\A}{\textbf{\textcolor{crimson}{Trả lời. }}}

% Khối "số liệu phải thuộc" / mẹo
\newtcolorbox{sohoc}{breakable,enhanced,sharp corners,
  boxrule=0.8pt,colframe=gold,colback=gold!6,
  coltitle=white,colbacktitle=gold,fonttitle=\bfseries\headfont\small,
  title={Số liệu \& mẹo phải thuộc},left=7pt,right=7pt,top=5pt,bottom=6pt}

\setlist[itemize]{leftmargin=1.4em,itemsep=1pt,topsep=2pt}

\begin{document}

% ─────────── Bìa ───────────
\begin{titlepage}
\centering
\vspace*{2.0cm}
{\headfont\bfseries\fontsize{15}{18}\selectfont\textcolor{soft}{ÔN VẤN ĐÁP BẢO VỆ ĐỒ ÁN — SỔ TAY CÁ NHÂN}}\\[6pt]
{\headfont\bfseries\fontsize{40}{44}\selectfont\textcolor{ink}{Đào Sỹ Duy Minh}}\\[6pt]
{\large\headfont\textcolor{crimson}{Project Manager \& AI/ML Engineer · MSSV 23122041}}\\[10pt]
{\color{crimson}\rule{9cm}{2pt}}\\[10pt]
{\large Dự án \textbf{StoryLens} — Nền tảng AI dịch \& hỏi--đáp Manga theo ngữ cảnh}\\[4pt]
{\color{soft}\large Trình bày slide 1--5 · phụ trách Q\&A: Quản lý dự án, AI/ML, chất lượng dịch, bản quyền}\\[2.2cm]
\begin{tabular}{rl}
\textbf{Môn học} & CSC10011 -- Công nghệ phần mềm cho hệ thống Trí tuệ Nhân tạo\\[3pt]
\textbf{GVHD} & GS.TS Nguyễn Văn Vũ, ThS. Trương Phước Lộc, ThS. Ngô Ngọc Đăng Khoa\\[3pt]
\textbf{Nhóm} & Nhóm 1 -- StoryLens · Ngày báo cáo 10/07/2026\\[3pt]
\textbf{Vai trò Q\&A} & PM (LN03) · Fine-tune OCR + Prompt dịch (LN07) · Bản quyền/Đạo đức (LN12)\\
\end{tabular}
\vfill
{\color{soft}\small Nhãn: \tagCore\ câu nền tảng phải thuộc · \tagTrap\ câu gài dễ mất điểm · \tagHonest\ phải trả lời trung thực, có định hướng}
\end{titlepage}

\tableofcontents
\newpage

% ═══════════════════════════════════════════════════════════════
\section{Chốt trong 30 giây: ``Minh đã làm gì?''}

\begin{qa}{Nếu hội đồng hỏi thẳng: bạn đóng góp cụ thể gì trong dự án?}
\A Em giữ \textbf{hai vai}: điều phối dự án và trực tiếp làm hai mắt xích AI quyết định chất lượng bản dịch.
\begin{itemize}
  \item \textbf{Quản lý (PM):} chọn quy trình \textbf{RUP rút gọn + Scrum} (sprint 2 tuần, 4 pha Inception\arr Elaboration\arr Construction\arr Transition); lập \& duy trì \textbf{Product/Sprint Backlog, WBS 31 task (135 story point), Gantt, Staff allocation, Risk register RE = P\x L, RTM, Burndown, Scrum log}; ước lượng bằng \textbf{Planning Poker + Story Points}; và bắt buộc \textbf{mọi Pull Request phải được review + CI xanh} mới merge vào \code{main}.
  \item \textbf{AI/ML (code thật):} \textbf{fine-tune manga-ocr} cho manga tiếng Nhật (\textbf{CER 6,1\%}, độ chính xác ký tự 93,9\%); \textbf{prompt engineering} cho \textbf{Gemini 2.5 Flash} dịch \textbf{theo cả trang} để giữ ngữ cảnh, xưng hô; và tối ưu quota API (caching, batch, xoay vòng key).
  \item \textbf{Task cụ thể trong WBS:} T1.2 (setup repo/board), \textbf{T2.1 fine-tune OCR trên Kaggle T4 (8 điểm)}, T4.1 (embedding pipeline / vector store), T5.3 (fix bug), T5.4 (tối ưu quota RPD), T6.5 (final deploy \& smoke test).
\end{itemize}
\textit{Câu nói mạnh:} ``Em vừa giữ nhịp dự án, vừa chịu trách nhiệm trực tiếp cho hai bước AI quyết định chất lượng: đọc chữ (OCR) và dịch có ngữ cảnh (prompt LLM).''
\end{qa}

% ═══════════════════════════════════════════════════════════════
\section{Phần Minh thuyết trình (slide 1--5) --- kể như một câu chuyện}

\begin{qa}{Slide 3 --- Vấn đề: vì sao thị trường cần StoryLens? \tagCore}
\A Cộng đồng đọc manga Việt (chủ yếu \textbf{15--28 tuổi}) rất lớn nhưng gặp \textbf{ba rào cản}:
\begin{itemize}
  \item \textbf{Fan-sub} ra chậm, chất lượng không đều, nhiều bộ bỏ dở vì thiếu nhân lực.
  \item Công cụ tự động (\textbf{Google Lens}) dịch \textbf{từng câu, mất ngữ cảnh} --- không giữ xưng hô/mạch truyện; DeepL/ChatGPT phải thao tác thủ công từng ảnh.
  \item Người đọc \textbf{không có cách hỏi ngay} về nội dung truyện.
\end{itemize}
Gói gọn: thị trường \textbf{thiếu một nền tảng dịch manga tự động vừa hiểu ngữ cảnh, vừa hỏi--đáp thông minh}.
\end{qa}

\begin{qa}{Slide 4--5 --- Định vị \& khác biệt (product position statement) \tagCore}
\A \textbf{Tuyên ngôn định vị (thuộc lòng):} ``Cho \textbf{người đọc manga Việt 15--28 tuổi} muốn đọc truyện Nhật chưa có bản dịch chính thức, StoryLens cho phép \textbf{tải trang lên, tự động dịch sang tiếng Việt có ngữ cảnh, và hỏi AI bất cứ điều gì về nội dung}.''
\begin{itemize}
  \item \textbf{Bốn giá trị cốt lõi:} (1) dịch giữ ngữ cảnh cả trang; (2) đọc overlay bật/tắt so sánh bản gốc; (3) hỏi--đáp \textbf{RAG có nguồn trích}; (4) hệ sinh thái cộng đồng (thư viện/share/forum).
  \item \textbf{Ma trận 4 tiêu chí:} dịch tự động · hiểu ngữ cảnh cả trang · RAG hỏi--đáp · overlay. Google Lens/DeepL/ChatGPT/fan-sub mỗi cái chỉ đạt 1--2 tiêu chí --- \textbf{chỉ StoryLens đủ cả bốn}.
\end{itemize}
\end{qa}

% ═══════════════════════════════════════════════════════════════
\section{Khối Quản lý dự án (LN03) --- Minh dễ bị hỏi nhất}

\begin{qa}{Vì sao chọn Agile/Scrum mà không dùng Waterfall? \tagCore}
\A Dự án \textbf{bất định công nghệ cao} (pipeline AI). Waterfall bắt chốt toàn bộ yêu cầu/kiến trúc từ đầu; nếu đến tuần 10 mới phát hiện OCR sai hoặc pipeline quá chậm thì sửa kiến trúc cực đắt, dễ sập tiến độ. Scrum cho phép \textbf{lặp \& tăng dần}: ngay từ Sprint 2 đã có pipeline AI thô chạy được để thử tích hợp \arr\ giảm rủi ro sớm, luôn có \textit{shippable increment} cuối mỗi sprint.
\end{qa}

\begin{qa}{Vì sao KHÔNG dùng COCOMO để ước lượng nỗ lực? \tagCore}
\A COCOMO cần ước lượng \textbf{KLOC} (nghìn dòng code) từ đầu --- với dự án AI/RAG, số dòng code \textbf{không phản ánh độ khó} (10 dòng prompt Gemini khó hơn hàng trăm dòng HTML). Nhóm dùng \textbf{Planning Poker + Story Points}, đánh giá theo \textit{độ phức tạp \x\ độ bất định (rủi ro) \x\ nỗ lực}. Mốc chuẩn: tính năng \textbf{Auth = 2 điểm}; \textbf{Pipeline AI = 8 điểm}. Các tính năng khác so với mốc để gán điểm.
\end{qa}

\begin{qa}{Velocity của nhóm là bao nhiêu? Tổng effort? \tagTrap}
\A \textbf{Trả lời đúng số thật trong file, KHÔNG nói ``20'':}
\begin{itemize}
  \item Tổng backlog \textbf{135 story point / 31 task}.
  \item \textbf{Velocity trung bình $\approx$ 22--23 điểm/sprint} (Burndown thật: 23 / 24 / 23 / 26 / 22 / 17 qua 6 sprint; trung bình 22,5).
\end{itemize}
\begin{center}\small
\begin{tabular}{lccc}
\toprule
Mốc & Điểm hoàn thành & Còn lại (actual) & \% \\
\midrule
End S1 & 23 & 112 & 17\% \\
End S2 & 24 & 88  & 35\% \\
End S3 & 23 & 65  & 52\% \\
End S4 & 26 & 39  & 71\% \\
End S5 & 22 & 17  & 87\% \\
End S6 & 17 & 0   & 100\% \\
\bottomrule
\end{tabular}
\end{center}
\textbf{Bẫy:} nếu nói cứng ``20 điểm'' rồi hội đồng mở sheet Burndown thấy 22,5 là lệch \arr\ mất uy tín số liệu.
\end{qa}

\begin{qa}{Là PM, bạn theo dõi tiến độ \& kiểm soát chất lượng thế nào? \tagCore}
\A Bốn cơ chế:
\begin{itemize}
  \item \textbf{Kanban} (To Do \arr\ In Progress \arr\ Reviewing \arr\ Done), thẻ chia nhỏ \lte\ 8 giờ.
  \item \textbf{Burndown chart} theo sprint (còn lại vs lý tưởng) để phát hiện chậm sớm.
  \item \textbf{Weekly Scrum} 3 câu hỏi (tuần qua làm gì / đang vướng gì / tuần tới làm gì) --- có \textbf{Scrum log 12 tuần} trong file kế hoạch.
  \item \textbf{Code review chéo:} không ai push thẳng \code{main}; mọi PR \gte\ 1 reviewer + CI xanh mới merge.
\end{itemize}
\end{qa}

\begin{qa}{Cho xem bằng chứng quản lý dự án thật (không phải nói miệng)? \tagCore}
\A Mở \code{docs/management/StoryLens\_Project\_Plan.xlsx} (10 sheet):
\begin{itemize}
  \item Sheet WBS: 31 task T1.1--T6.5, người phụ trách, output, điểm, trạng thái.
  \item Sheet \textbf{Risk Register}: RE = Probability \x\ Loss số hoá theo bảng LN03, có rank.
  \item Sheet Gantt (task \x\ 12 tuần + mốc Build), Staff Allocation (6 người \x\ 12 tuần), RTM, Burndown, Scrum log.
\end{itemize}
\end{qa}

\begin{qa}{Rủi ro lớn nhất của dự án và cách ứng phó? (Risk register) \tagCore}
\A RE = Probability \x\ Loss (thang LN03). Top rủi ro:
\begin{center}\small
\begin{tabular}{clcc p{4.6cm}}
\toprule
ID & Rủi ro & RE & Rank & Ứng phó \\
\midrule
R2 & Quota/chi phí LLM API vượt budget & 5625 & 1 & Cache; key-pool 4 key ($\sim$2000 RPD); fallback free tier \\
R5 & Dữ liệu manga vướng bản quyền & 5625 & 2 & Dùng dataset public (Manga109-s); nội dung do user đăng \\
R1 & OCR accuracy $<$ 95\% (CER cao) & 3750 & 3 & Thêm dữ liệu; fine-tune manga-ocr \\
R4 & ChromaDB không scale & 3750 & 4 & \textbf{Đã chuyển sang pgvector} \\
\bottomrule
\end{tabular}
\end{center}
\end{qa}

\begin{qa}{Bạn là PM nhưng slide nói ``flat team không phân cấp'' --- có mâu thuẫn? \tagTrap}
\A Không. \textbf{Flat team} nói về việc \textit{quyết định kỹ thuật được thảo luận tập thể}, không có cấp trên áp đặt giải pháp. Vai trò PM của em là \textbf{điều phối tiến độ, quản lý backlog, quản trị rủi ro, và giữ cổng chất lượng (PR review + CI)} --- là trách nhiệm quy trình, không phải quyền lực ra lệnh kỹ thuật. Hai điều này không xung đột.
\end{qa}

% ═══════════════════════════════════════════════════════════════
\section{Khối AI/ML của Minh (LN07) --- OCR + prompt dịch}

\begin{qa}{Nêu 8 bước ML Workflow nhóm áp dụng. \tagCore}
\A Requirements \arr\ Data collection \& preparation \arr\ Feature engineering \arr\ Model training \arr\ \textbf{Model evaluation} \arr\ Model deployment \arr\ Operating \& monitoring \arr\ Maintenance (retrain từ ca lỗi khó: chữ SFX, viết tay).
\end{qa}

\begin{qa}{(BẪY LỚN NHẤT) Bạn fine-tune manga-ocr thế nào? Chứng minh nó tốt hơn model gốc? \tagTrap}
\A \textbf{Có ablation hợp lệ + artifact xác minh} --- số liệu thật trong \code{model\_evaluation.docx} (PA4):
\begin{itemize}
  \item \textbf{Cách làm:} base \code{kha-white/manga-ocr-base} (VisionEncoderDecoderModel), \textbf{fine-tune full} trên \textbf{Manga109-s}. Dataset OCR \textbf{123.208 cặp} (crop bong bóng + nhãn chữ Nhật), chia \textbf{split-by-title, seed = 42} (chống rò rỉ) \arr\ Train 97.602 / Val 11.476 / \textbf{Test 14.130}.
  \item \textbf{Cấu hình:} 8 epochs, AdamW, lr = 3e-5, cosine schedule, warmup 5\%, label-smoothing = 0.1, weight\_decay = 0.01, fp16; chọn best theo \textbf{CER val nhỏ nhất}, early-stop patience 3; train trên \textbf{Kaggle Tesla T4}.
  \item \textbf{Bằng chứng fine-tune hiệu quả (ablation 2 vòng, cùng dataset/test-split/base):}
\end{itemize}
\begin{center}\small
\begin{tabular}{lcc}
\toprule
Phiên bản & Char-Acc & CER \\
\midrule
PA2 (fine-tune sơ bộ, ít epoch) & 87,3\% & \textbf{12,7\%} \\
PA3 (fine-tune hoàn chỉnh) & 93,9\% & \textbf{6,1\%} \\
\bottomrule
\end{tabular}
\end{center}
\begin{itemize}
  \item \textbf{CER giảm 12,7\% \arr\ 6,1\% ($\approx$ giảm hơn 50\% lỗi)}. Khớp \textbf{chính xác 4 chữ số} với artifact thật \code{ai\_module/models/manga\_ocr/weights/metadata.json} (test\_cer = 0.061090) \arr\ chứng minh là output huấn luyện thật, không phải số gõ tay.
\end{itemize}
\end{qa}

\begin{qa}{Demo sống n=1 cho thấy fine-tune đọc SAI chữ đầu --- vậy nó tệ hơn base? \tagTrap}
\A Đó là \textbf{1 mẫu (n=1) trên ảnh minh hoạ AI-generated ngoài phân phối} Manga109-s --- chỉ là smoke-test định tính, \textbf{không dùng để kết luận}. Bằng chứng chuẩn là \textbf{test split 14.130 mẫu}: CER 12,7\% \arr\ 6,1\%. Kết quả n=1 chỉ phản ánh: mô hình được tối ưu cho \textbf{phong cách line-art Manga109-s}, chưa tổng quát hoá hoàn hảo sang mọi phong cách vẽ --- đúng như giới hạn đã ghi nhận trong SAD.
\end{qa}

\begin{qa}{Giải thích công thức \& ý nghĩa CER, mAP@0.5, IoU, Precision. \tagCore}
\A
\begin{itemize}
  \item \textbf{CER} (đánh giá OCR --- phần của Minh): $CER = \dfrac{S+D+I}{N}$ ($S$ thay thế sai, $D$ xoá/bỏ sót, $I$ chèn thừa, $N$ tổng ký tự chuẩn). \textbf{6,1\%} = trong 100 chữ chỉ sai $\sim$6.
  \item \textbf{IoU} = diện tích giao / diện tích hợp giữa khung dự đoán \& khung thật; ngưỡng \gte\ 0.5 tính là đúng.
  \item \textbf{Precision} $= \dfrac{TP}{TP+FP} = \textbf{97,3\%}$ --- rất ít vẽ nhầm khung dịch vào chi tiết phong cảnh.
  \item \textbf{mAP@0.5} = trung bình Average Precision (diện tích dưới PR-curve) tại IoU \gte\ 0.5 $= \textbf{95,3\%}$; Recall 93,8\%; mAP@0.5:0.95 = 84,4\%.
  \item \textbf{Dữ liệu đánh giá:} Manga109-s, test \textbf{956 ảnh / $>$14.130 vùng chữ}, ngưỡng đặt 85\% \arr\ đều vượt.
\end{itemize}
\end{qa}

\begin{qa}{Vì sao chọn YOLOv8n (nano) chứ không model lớn hơn cho chính xác hơn? \tagCore}
\A YOLOv8n chỉ \textbf{$\sim$3M tham số}, chạy được \textbf{CPU-only} trên HuggingFace Spaces free tier (không GPU), độ trễ detection $\sim$\textbf{0,13 giây/ảnh}; vẫn đạt \textbf{mAP@0.5 = 95,3\% $>$ ngưỡng 85\%}. Chọn nano là cân bằng chính xác/chi phí hợp lý cho hạ tầng miễn phí --- dùng model lớn hơn sẽ tràn RAM và chậm mà lợi ích biên không đáng kể.
\end{qa}

\begin{qa}{Prompt dịch của bạn cụ thể là gì? ``Context injection'' nghĩa là gì? \tagCore}
\A Module dịch (task T3.2: \code{GlossaryManager} + \code{LLMClient} gọi Gemini) đưa \textbf{cả trang + glossary xưng hô} vào prompt thay vì dịch từng bong bóng rời. Nhờ ``\textbf{context injection}'' này, Gemini giữ \textbf{nhất quán xưng hô, mạch truyện, sắc thái} trên toàn trang. Kèm \textbf{xoay vòng nhiều API key} chống hết quota và \textbf{fallback M2M100/NLLB} khi Gemini lỗi.
\end{qa}

\begin{qa}{OCR sai thì Gemini tự sửa --- nhưng nếu OCR sai nặng, Gemini bịa (hallucinate) thì sao? \tagTrap}
\A Đúng, đây là rủi ro ``Artificial Stupidity''. Nhóm giảm bằng ba lớp: (1) dịch \textbf{cả trang có ngữ cảnh} nên Gemini có nhiều tín hiệu để suy đoán đúng; (2) \textbf{cảnh báo khi độ tin cậy OCR thấp} ngay trên reader; (3) \textbf{Studio QC} cho dịch giả sửa tay từng bong bóng, lưu lịch sử hiệu đính. \textbf{AI không có tiếng nói cuối cùng --- con người mới quyết định} (human-in-the-loop).
\end{qa}

\begin{qa}{CER 6,1\% đo trên Manga109-s. User upload manga phong cách khác thì CER thật bao nhiêu? \tagHonest}
\A Trung thực: 6,1\% là \textbf{in-distribution} (đúng phân phối huấn luyện). Out-of-distribution \textbf{chưa đo bằng số} (mới có 1 mẫu định tính) --- đây là giới hạn đã ghi rõ trong báo cáo. Nhóm bù bằng \textbf{Gemini self-correction} và \textbf{Studio QC}. Hướng phát triển: chạy đánh giá trên vài chục mẫu chữ Nhật ngoài phân phối để có ước lượng CER thô.
\end{qa}

% ═══════════════════════════════════════════════════════════════
\section{Các câu Q\&A Minh được phân công (từ phụ lục kịch bản)}

\begin{qa}{Làm sao đảm bảo chất lượng bản dịch? (Minh / Phú Hòa) \tagCore}
\A Gemini dịch \textbf{theo cả trang} nên giữ ngữ cảnh, xưng hô; cộng \textbf{human-in-the-loop}: Studio QC cho duyệt/sửa từng bong bóng và cơ chế phản hồi tốt/xấu (endpoint \code{/v1/pages/\{id\}/feedback}) để cải thiện dần. Ngưỡng chấp nhận theo SAD: tỉ lệ 👍 \gte\ 80\%.
\end{qa}

\begin{qa}{Vấn đề bản quyền manga thì sao? (Minh / Quý) \tagHonest}
\A StoryLens là \textbf{công cụ hỗ trợ đọc/dịch}, không phải nơi phát tán. Ba lớp: (1) mô hình huấn luyện trên \textbf{dataset học thuật Manga109-s} (giấy phép nghiên cứu); (2) nội dung người dùng upload mặc định \textbf{Private}, có \textbf{điều khoản sử dụng + công cụ kiểm duyệt admin}; (3) roadmap có \textbf{quy trình DMCA take-down}. Đã ghi nhận là \textbf{Risk R5} trong risk register.
\end{qa}

\begin{qa}{Điểm khác biệt lớn nhất so với ChatGPT / Google Lens? (Minh) \tagCore}
\A \textbf{Ba-trong-một chuyên biệt} mà công cụ chung không có: (1) dịch \textbf{cả trang giữ ngữ cảnh}; (2) \textbf{render overlay đúng bong bóng} (binary-search chọn cỡ chữ vừa khít); (3) \textbf{RAG hỏi--đáp giới hạn theo đúng thư viện riêng của từng người dùng, có nguồn trích}. ChatGPT/Lens chỉ làm một phần, không render vào bong bóng, không hỏi--đáp trên đúng nội dung truyện của bạn.
\end{qa}

\begin{qa}{Vì sao tài liệu PA ghi ChromaDB / all-MiniLM / GPT-4 mà sản phẩm là pgvector / Gemini? (Phú Hòa / Minh) \tagTrap}
\A Đây là \textbf{sự tiến hoá qua các vòng lặp RUP} --- điểm mạnh về kỹ năng ra quyết định, không phải lỗi:
\begin{itemize}
  \item Kế hoạch ban đầu: ChromaDB + embedding all-MiniLM + GPT-4.
  \item Khi hiện thực: hợp nhất về \textbf{pgvector ngay trong Supabase} (bớt một dịch vụ, đỡ vận hành --- \textbf{Risk R4} ghi rõ ``ChromaDB không scale \arr\ đã đổi pgvector''); chọn \textbf{Gemini} để đồng bộ (vừa dịch, vừa embedding \textbf{768 chiều}, vừa RAG, xoay vòng key rẻ).
  \item \textbf{Sản phẩm cuối phản ánh quyết định kỹ thuật tốt hơn bản kế hoạch.}
\end{itemize}
\end{qa}

\begin{qa}{Đóng góp từng thành viên / cách làm việc nhóm? (Minh) \tagCore}
\A \textbf{Flat team 6 người}, RUP rút gọn + Scrum, sprint 2 tuần, weekly scrum, mọi PR review:
\begin{itemize}
  \item \textbf{Minh (23122041)} --- PM \& AI/ML: quản lý tiến độ, fine-tune OCR, prompt dịch.
  \item \textbf{Hà Dương (23122002)} --- AI/ML (OCR): dataset, hỗ trợ fine-tune, đánh giá accuracy.
  \item \textbf{Kiệt (23122039)} --- Backend: FastAPI, DB schema, tích hợp pipeline.
  \item \textbf{Phú Hòa (23122030)} --- AI/ML (RAG): Gemini, pgvector, pipeline hỏi--đáp.
  \item \textbf{Nguyên (23122044)} --- Frontend: Next.js Upload/Reader/Q\&A.
  \item \textbf{Phú Quý (23122048)} --- QA \& BA: SDP/Vision/Use case, test case, UAT.
\end{itemize}
\end{qa}

% ═══════════════════════════════════════════════════════════════
\section{Bẫy mâu thuẫn tài liệu (câu ``bắt lỗi'')}

\begin{qa}{WBS ghi task ChromaDB/FAISS (T2.3, T4.1) nhưng sản phẩm dùng pgvector --- giải thích? \tagTrap}
\A WBS là \textbf{kế hoạch ở thời điểm lập}; kiến trúc tiến hoá qua các sprint là chuyện bình thường trong RUP/Agile. Bằng chứng nhóm \textit{chủ động} ra quyết định nằm ngay trong \textbf{Risk register R4} (``ChromaDB không scale \arr\ chuyển pgvector''). Không phải sai lệch tài liệu mà là \textbf{quản trị thay đổi có ghi vết}.
\end{qa}

\begin{qa}{Vì sao chọn embedding 768 chiều? Có phải tuỳ tiện? \tagCore}
\A 768 chiều là kích thước vector chuẩn của \code{gemini-embedding-001}. Chọn nó để \textbf{đồng bộ toàn hệ Gemini} (một nhà cung cấp cho cả dịch, embedding, RAG \arr\ đơn giản vận hành, chung cơ chế xoay vòng key), và pgvector index thẳng vector 768 chiều với cosine similarity.
\end{qa}

% ═══════════════════════════════════════════════════════════════
\section{Đạo đức AI \& bản quyền (LN12) --- vùng của Minh}

\begin{qa}{Train trên Manga109-s và cho user dịch truyện có bản quyền --- có vi phạm? \tagHonest}
\A Tách hai chuyện: (1) \textbf{huấn luyện} dùng \textbf{Manga109-s} --- dataset học thuật có giấy phép nghiên cứu, hợp lệ cho mục đích nghiên cứu/giáo dục; (2) \textbf{nội dung người dùng upload} do người dùng tự chịu trách nhiệm, mặc định \textbf{Private}, có điều khoản + kiểm duyệt admin + roadmap DMCA. Đây là \textbf{Risk R5} đã nhận diện và có chiến lược ứng phó.
\end{qa}

\begin{qa}{AI Bias trong hệ thống là gì? \tagHonest}
\A YOLOv8 \& OCR huấn luyện chủ yếu trên \textbf{Manga109-s (manga Nhật)} \arr\ \textbf{distribution bias}: giỏi với phong cách Nhật (bong bóng tròn, line-art), \textbf{kém hơn với manhwa Hàn / manhua Trung}. Nhóm \textbf{ghi nhận rõ giới hạn} trong tài liệu; hướng khắc phục = bổ sung dataset đa dạng để giảm bias. Theo LN12: ``Principles alone cannot guarantee ethical AI'' --- nên nhóm \textbf{đo và công bố giới hạn}, không chỉ hứa ``cố gắng công bằng''.
\end{qa}

\begin{qa}{Hệ thống có theo dõi chất lượng mô hình AI sau deploy không (MLOps)? \tagHonest}
\A Có đủ \textbf{giám sát hạ tầng} (Sentry bắt lỗi, OpenTelemetry theo dõi latency, Admin health dashboard, \code{/health/deep}). \textbf{Model monitoring} (theo dõi drift mAP/CER theo thời gian) \textbf{chưa có} --- vì đang ở giai đoạn MVP, model chạy tĩnh (update thủ công). Hướng phát triển: tích hợp \textbf{W\&B / MLflow} để theo dõi drift và trigger retrain tự động.
\end{qa}

% ═══════════════════════════════════════════════════════════════
\section{15 câu bẫy khó nhất --- luyện phản xạ nhanh}

\begin{enumerate}[leftmargin=1.6em,itemsep=3pt]
  \item \textbf{Fine-tune OCR thế nào?} \arr\ base \code{manga-ocr-base}, fine-tune full Manga109-s (123k cặp, split-by-title seed 42), 8 epochs AdamW lr=3e-5 cosine, T4 Kaggle. CER 12,7\% \arr\ 6,1\%, xác minh bằng \code{metadata.json}.
  \item \textbf{Demo n=1 fine-tune đọc sai --- tệ hơn base?} \arr\ n=1 ảnh ngoài phân phối, smoke-test; chuẩn là test split 14.130 mẫu.
  \item \textbf{CER thực khi manga khác phong cách?} \arr\ Trung thực: chỉ đo in-distribution; OOD chưa đo; bù bằng Gemini + Studio QC.
  \item \textbf{OCR sai \arr\ Gemini bịa?} \arr\ Rủi ro có thật; giảm bằng dịch cả trang + cảnh báo confidence + human-in-the-loop.
  \item \textbf{Prompt dịch \& context injection?} \arr\ \code{GlossaryManager}+\code{LLMClient}, đưa cả trang + glossary xưng hô vào prompt; fallback M2M100/NLLB.
  \item \textbf{Vì sao YOLOv8n nano?} \arr\ 3M tham số, CPU-only free tier, 0,13s/ảnh, vẫn mAP 95,3\% $>$ 85\%.
  \item \textbf{Velocity? Tổng effort?} \arr\ $\sim$22--23 SP/sprint, tổng 135 SP / 31 task. (KHÔNG nói ``20''.)
  \item \textbf{Vì sao không COCOMO?} \arr\ Cần KLOC, không đo được độ khó AI; dùng Planning Poker + Story Points.
  \item \textbf{PM mà flat team --- mâu thuẫn?} \arr\ Flat = quyết định kỹ thuật tập thể; PM lo tiến độ/backlog/rủi ro/PR gate.
  \item \textbf{Bằng chứng PM thật?} \arr\ \code{StoryLens\_Project\_Plan.xlsx}: WBS, Gantt, Risk RE=P\x L, RTM, Burndown, Scrum log.
  \item \textbf{ChromaDB trong WBS vs pgvector?} \arr\ Tiến hoá RUP; Risk R4 ghi rõ ``đã đổi pgvector''.
  \item \textbf{GPT-4/all-MiniLM vs Gemini?} \arr\ Cùng lý do tiến hoá; Gemini đồng bộ dịch+embedding 768+RAG, xoay vòng key rẻ.
  \item \textbf{Bản quyền?} \arr\ Train Manga109-s (license nghiên cứu); user content Private + kiểm duyệt + DMCA roadmap; Risk R5.
  \item \textbf{Khác ChatGPT/Lens?} \arr\ 3-trong-1: dịch cả trang + overlay render + RAG scoped có nguồn trích.
  \item \textbf{Model monitoring sau deploy?} \arr\ Có infra monitoring (Sentry/OTel/health); model drift chưa có, roadmap W\&B/MLflow.
\end{enumerate}

% ═══════════════════════════════════════════════════════════════
\section{Phụ lục --- Số liệu \& điểm yếu phải thủ sẵn}

\begin{sohoc}
\textbf{Số phải bật ra ngay khi bị hỏi:}
\begin{itemize}
  \item OCR: \textbf{CER 6,1\%} (Char-Acc 93,9\%, exact-match 67,5\%); ablation \textbf{12,7\% \arr\ 6,1\%}.
  \item Detection: \textbf{mAP@0.5 = 95,3\%}, mAP@0.5:0.95 = 84,4\%, Precision 97,3\%, Recall 93,8\% (ngưỡng 85\%).
  \item Dataset: Manga109-s, OCR \textbf{123.208 cặp} (80/10/10, seed 42); detection test \textbf{956 ảnh / 14.130 vùng chữ}.
  \item PM: \textbf{135 SP / 31 task}, velocity \textbf{$\sim$22--23 SP/sprint}, 6 sprint \x\ 2 tuần, 3 mốc Build.
  \item Kiến trúc: \textbf{3 dịch vụ} (Vercel / Render 512MB / HF Spaces $\sim$4GB), embedding \textbf{768 chiều}, CI \textbf{5 job}.
  \item Rủi ro top: R2 quota \& R5 bản quyền (RE = 5625).
\end{itemize}
\end{sohoc}

\begin{tcolorbox}[breakable,enhanced,sharp corners,boxrule=0.8pt,colframe=leaf,colback=leaf!5,
  coltitle=white,colbacktitle=leaf,fonttitle=\bfseries\headfont\small,
  title={Điểm yếu phải thừa nhận trung thực (kèm định hướng) --- đừng để reviewer bắt trước}]
\begin{itemize}
  \item \textbf{YOLO thiếu artifact gốc} (chỉ có \code{best.pt}, không có \code{results.csv}/PR-curve) \arr\ số mAP dựa vào SAD PA3; hướng khắc phục: lưu \code{results.csv} khi train lại.
  \item \textbf{CER out-of-distribution chưa đo} (mới n=1 định tính) \arr\ hướng: đánh giá vài chục mẫu OOD.
  \item \textbf{Không có Model Monitoring} (drift) \arr\ đang MVP, roadmap W\&B/MLflow.
  \item \textbf{Defect log rỗng} (61/61 Pass, 0 defect) \arr\ nếu bị nghi ``test cho có'', trích 1--2 bug thật từ git history.
  \item \textbf{Chưa có dữ liệu feedback dịch sống} (không truy cập Supabase production) \arr\ mô tả đúng cơ chế đã hiện thực, không suy diễn số.
\end{itemize}
\end{tcolorbox}

\vspace{6pt}
\begin{center}\small\color{soft}
Sổ tay cá nhân --- Đào Sỹ Duy Minh · StoryLens · CSC10011 (HCMUS) · in A4 cầm theo khi bảo vệ.
\end{center}

\end{document}
